\documentclass[11pt,a4paper]{article}
\usepackage[utf8]{inputenc}
\usepackage[T1]{fontenc}
\usepackage[french]{babel}
\frenchbsetup{StandardLists=true}
\usepackage[left=2cm,right=2cm,top=2cm,bottom=2cm]{geometry}
\usepackage[dvipsnames]{xcolor}
\usepackage{fouriernc}
\usepackage{amsmath,amssymb,amsfonts}
\usepackage{array,tabularx,booktabs,multirow}
\usepackage{colortbl}
\usepackage{graphicx}
\usepackage{mdframed}
\usepackage{enumitem}
\usepackage{fancyhdr}
\usepackage{chemformula}
\usepackage{awesomebox}
\setlength{\aweboxvskip}{0mm}
\usepackage{tcolorbox}
\tcbuselibrary{skins}
\usepackage{fontawesome5}

\definecolor{exoheader}{HTML}{1E5AA0}
\definecolor{exolight}{HTML}{DCE8FF}
\definecolor{warnorange}{RGB}{220,100,0}
\definecolor{problemebg}{HTML}{FFFDE7}

\renewcommand{\arraystretch}{1.8}
\setlength{\extrarowheight}{1mm}

\pagestyle{fancy}
\renewcommand\headrulewidth{1pt}
\fancyhead[L]{Académie de Besançon}
\fancyhead[C]{\textbf{TP --- Extraction et CCM de la vanilline}}
\fancyhead[R]{Terminale / BTS Chimie}
\fancyfoot[C]{page \thepage}
\fancyfoot[R]{2025 -- 2026}
\fancyfoot[L]{Alexis RUIZ}

\newcommand{\sectitle}[1]{%
  \vspace{6pt}
  \noindent\colorbox{exoheader}{\parbox{\dimexpr\linewidth-2\fboxsep}%
    {\color{white}\bfseries\large\strut\quad #1}}
  \vspace{4pt}
}
\newcommand{\subsectitle}[1]{%
  \vspace{4pt}
  \noindent\colorbox{exolight}{\parbox{\dimexpr\linewidth-2\fboxsep}%
    {\color{exoheader}\bfseries\strut\quad #1}}
  \vspace{2pt}
}
\newcounter{question}
\newcommand{\question}{%
  \stepcounter{question}%
  \noindent\textcolor{exoheader}{\textbf{Question \thequestion.}}\quad
}

\begin{document}

\noindent\textbf{NOM :}\hfill\textbf{Prénom :}\hfill\textbf{Classe :}\hfill\phantom{x}\\[0.2cm]
\hrule\vspace{0.2cm}

\begin{center}
  {\large\bfseries\textcolor{exoheader}{Travaux Pratiques de Chimie Organique}}\\[4pt]
  {\LARGE\bfseries Extraction et identification de la vanilline (CCM)}\\[2pt]
  {\small Académie de Besançon --- Belfort}
\end{center}
\hrule\vspace{6pt}

\begin{tcolorbox}[enhanced, colback=problemebg, colframe=exoheader, boxrule=1.5pt,
  title={\faFlask\quad Problématique},
  fonttitle=\bfseries\color{white},
  attach boxed title to top left={yshift=-2mm, xshift=4mm},
  boxed title style={colback=exoheader, colframe=exoheader}]
\textit{De nombreux produits alimentaires sont parfumés à la vanille.
Certains contiennent de la \textbf{vraie vanilline} (extraite de gousses), d'autres de la
\textbf{vanilline de synthèse}. \textbf{Comment identifier et distinguer ces produits ?}}
\end{tcolorbox}

\vspace{6pt}

% ══════════════════════════════════════════════════════════════════════════════
\sectitle{Partie I --- Extraction de l'arôme par solvant}
% ══════════════════════════════════════════════════════════════════════════════

\begin{mdframed}[backgroundcolor=exolight, linecolor=exoheader, linewidth=1.2pt]
L'extraction par solvant consiste à transférer un soluté d'une phase aqueuse vers
une phase organique non miscible (ici l'éthanoate d'éthyle).
\end{mdframed}

\vspace{4pt}
\question Pourquoi utilise-t-on l'éthanoate d'éthyle comme solvant d'extraction ?
Quelle propriété doit-il avoir vis-à-vis de l'eau ? \dotfill

\vspace{0.3cm}\noindent\dotfill

\vspace{4pt}
\subsectitle{A --- Extraction du sucre vanillé (solution 1)}

\begin{enumerate}[leftmargin=1.5cm, label=\textcolor{exoheader}{\textbf{\arabic*.}}]
  \item Peser environ 5~g de sucre vanillé dans un erlenmeyer.
  \item Dissoudre dans 30~mL d'eau distillée (agiter).
  \item Ajouter 10~mL d'éthanoate d'éthyle et une spatulée de \ch{NaCl}.
  \item Verser dans l'ampoule à décanter ; agiter en dégazant régulièrement.
  \item Laisser décanter, éliminer la phase inférieure (aqueuse) dans un bécher.
  \item Recueillir la phase supérieure (organique) dans un flacon étiqueté \og Solution 1 \fg{}.
\end{enumerate}

\vspace{4pt}
\subsectitle{B --- Extraction du sucre vanilliné (solution 2)}

Même protocole avec 5~g de sucre vanilliné. Recueillir dans un flacon étiqueté \og Solution 2 \fg{}.

\vspace{4pt}
\subsectitle{C --- Extraction d'une gousse de vanille (solution 3)}

\begin{enumerate}[leftmargin=1.5cm, label=\textcolor{exoheader}{\textbf{\arabic*.}}]
  \item Peser environ 1~g de gousse de vanille, la découper en petits morceaux.
  \item Introduire dans un erlenmeyer avec 10~mL d'éthanoate d'éthyle.
  \item Agiter 20~minutes à l'agitateur magnétique, puis filtrer.
  \item Recueillir le filtrat dans un flacon étiqueté \og Solution 3 \fg{}.
\end{enumerate}

\vspace{4pt}
\question Dans quelle phase se trouve la vanilline après extraction ? Justifier. \dotfill

\vspace{0.3cm}\noindent\dotfill

% ══════════════════════════════════════════════════════════════════════════════
\newpage
\sectitle{Partie II --- Identification par chromatographie (CCM)}
% ══════════════════════════════════════════════════════════════════════════════

\begin{mdframed}[backgroundcolor=exolight, linecolor=exoheader, linewidth=1.2pt]
La \textbf{chromatographie sur couche mince (CCM)} permet de séparer et d'identifier
des espèces chimiques en comparant leur rapport frontal $R_f$ avec celui d'un témoin.
\[
  R_f = \frac{\text{distance parcourue par le soluté}}{\text{distance parcourue par l'éluant}}
\]
\end{mdframed}

\vspace{4pt}
\subsectitle{1. Préparation de la cuve}

\begin{enumerate}[leftmargin=1.5cm, label=\textcolor{exoheader}{\textbf{\arabic*.}}]
  \item Mesurer 10~mL d'éluant (cyclohexane + éthanoate d'éthyle) et le verser dans la cuve.
  \item Fermer la cuve et laisser saturer l'atmosphère quelques minutes.
\end{enumerate}

\vspace{4pt}
\subsectitle{2. Dépôt des solutions}

\begin{enumerate}[leftmargin=1.5cm, label=\textcolor{exoheader}{\textbf{\arabic*.}}]
  \item Sur une plaque CCM ($5 \times 10$~cm, gel de silice), tracer au crayon une ligne de dépôt à 1~cm du bas.
  \item Repérer 4 points (1 cm d'écart minimum) : sol. 1, sol. 2, sol. 3, vanilline (témoin, sol. 4).
  \item Déposer à chaque point 2 à 3 gouttes avec un capillaire ; sécher entre chaque dépôt.
\end{enumerate}

\vspace{4pt}
\subsectitle{3. Élution et révélation}

\begin{enumerate}[leftmargin=1.5cm, label=\textcolor{exoheader}{\textbf{\arabic*.}}]
  \item Placer délicatement la plaque dans la cuve (dépôts au-dessus du solvant).
  \item Laisser migrer jusqu'à 1~cm du bord supérieur, retirer et marquer le front du solvant.
  \item Révéler sous lampe UV (254~nm) et entourer les taches au crayon.
\end{enumerate}

\vspace{4pt}
\question Reproduire ci-dessous le schéma de la plaque obtenue et indiquer les taches observées.

\vspace{0.5cm}
\begin{center}
\begin{tikzpicture}[scale=0.9]
\draw[thick] (0,0) rectangle (8,6);
\draw[dashed] (0,0.6) -- (8,0.6) node[right]{\small ligne de dépôt};
\draw[dashed] (0,5.4) -- (8,5.4) node[right]{\small front du solvant};
\foreach \x/\lab in {1/sol.1, 2.5/sol.2, 4/sol.3, 5.5/témoin}{
  \node[below] at (\x,0.6) {\small\lab};
}
\end{tikzpicture}
\end{center}

\vspace{4pt}
\question Calculer le $R_f$ de la vanilline (témoin) et comparer avec les autres solutions.

\vspace{1.5cm}

\question Lequel ou lesquels des produits testés contient de la vanilline ? Justifier. \dotfill

\vspace{0.3cm}\noindent\dotfill

\vspace{4pt}
\question Le sucre vanilliné contient-il la même espèce que le sucre vanillé ?
Quelles conclusions pouvez-vous en tirer ? \dotfill

\vspace{0.3cm}\noindent\dotfill

\vspace{0.3cm}\noindent\dotfill

\end{document}
